\begin{trapbox}
\textbf{\textcolor{CTrap}{易错一｜把曲线外点错当切点}}

\textbf{错误做法：} 只要给出 $Q(x_1,y_1)$，就写
$\dfrac{x_1x}{a^2}+\dfrac{y_1y}{b^2}=1$ 并称其为“在 $Q$ 处的切线”。

\textbf{纠正：} 首先检查 $\dfrac{x_1^2}{a^2}+\dfrac{y_1^2}{b^2}=1$。
若 $Q$ 在椭圆外，该式其实是 $Q$ 对应的\textbf{切点弦（极线）}方程，通常连结从 $Q$ 引出的两条切线的切点，\textbf{不是}过 $Q$ 的切线。

\medskip
\textbf{\textcolor{CTrap}{易错二｜端点处分母为零}}

\textbf{错误做法：} 当 $y_0=0$，继续套
$k=-\dfrac{b^2x_0}{a^2y_0}$。

\textbf{纠正：} 若 $y_0=0$，$P=(\pm a,0)$，竖直切线 $x=\pm a$；
若 $x_0=0$，$P=(0,\pm b)$，水平切线 $y=\pm b$。
统一公式 $\dfrac{x_0x}{a^2}+\dfrac{y_0y}{b^2}=1$ 在两种情况下都可以直接使用。

\medskip
\textbf{\textcolor{CTrap}{易错三｜把交换因子次序当成错误}}

$\dfrac{x_0x}{a^2}$ 与 $\dfrac{xx_0}{a^2}$ \textbf{完全相同}，不是两种不同公式。
真正容易错的是：把分母 $a^2,b^2$ 交换、把椭圆的加号误写成减号，或者把右侧 $1$ 写错。

\medskip
\textbf{\textcolor{CTrap}{易错四｜把圆的半径直接当椭圆法线}}

圆 $x^2+y^2=r^2$ 在 $P$ 处法向量为 $(x_0,y_0)$；
椭圆 $x^2/a^2+y^2/b^2=1$ 的法向量却是
$\left(x_0/a^2,y_0/b^2\right)$。
只有特殊位置（如轴端点）或圆的情形，$OP$ 才可能恰好沿法线方向。

\medskip
\textbf{\textcolor{CTrap}{易错五｜只套 $y=kx+m$，漏掉竖直切线}}

非竖直切线满足 $m^2=a^2k^2+b^2$，但\textbf{竖直切线} $x=\pm a$ 无法写成 $y=kx+m$，必须单列。
此外，$m=0$ 不可能满足上述切线判定，因为 $a^2k^2+b^2>0$。

\medskip
\textbf{\textcolor{CTrap}{易错六｜截距式除以零、面积漏绝对值}}

只有 $x_0y_0\ne0$ 时才可写
$\dfrac{x}{a^2/x_0}+\dfrac{y}{b^2/y_0}=1$，并计算有限的两轴截距围成的三角形面积
\[
 S=\frac12\left|\frac{a^2}{x_0}\frac{b^2}{y_0}\right|
 =\frac{a^2b^2}{2|x_0y_0|}.
\]
当切线与一条坐标轴平行时，没有通常意义下的有限截距三角形，不能套此面积公式。

\medskip
\textbf{\textcolor{CTrap}{易错七｜把“平方替换法”机械套到一切二次曲线}}

对椭圆、双曲线的标准方程，平方项分别替换为相应的乘积确实有效；
对抛物线 $y^2=2px$，正确切线是
$y_0y=p(x+x_0)$，\textbf{不是} $y_0y=2px$。

\medskip
\textbf{\textcolor{CTrap}{自检三问}}

\textbf{点真在椭圆上吗？}\quad
\textbf{是已知切点，还是已知外点／斜率？}\quad
\textbf{是否遗漏竖直切线？}
\end{trapbox}
